% GENERATED FILE. Edit canonical lessons, then run npm run generate:books.
% canonical-source-hash: fnv1a64:134c6f6050685816
% canonical-lessons: IT-C269-ero-contento, IT-C269-spero-di, IT-C269-ripasso-sentimenti

\chapter{How It Felt --- Feelings and Hopes}
\label{ch:it-feelings-and-hopes}
\hlchaptermodality{269}

\begin{chapteropening}
\textbf{By the end of this chapter:} \emph{I can say how an experience felt, in the past, and what I hope for, with spero di and il mio sogno è.}

Complete IT-C269-ripasso-sentimenti: a journey told from both ends, the feeling looking back and the hope looking forward.
\end{chapteropening}

\section[ero contento, ero contenta]{ero contento — how it felt}

\label{lesson:IT-C269-ero-contento}



\begin{quote}
Say \textquotedblleft{}I was tired, but I was happy\textquotedblright{}. (\emph{Ero stanco, ma ero contento}, or \emph{ero stanca, ma ero contenta}.)

You said that sentence in chapter two hundred and sixty-three. This lesson is about why it worked.
\end{quote}

\begin{grammarlens}[title={Grammar Lens: a feeling is scene}]
You already own the feelings: \textbf{contento}, \textbf{felice}, \textbf{triste}, \textbf{stanco}, \textbf{nervoso}. Describing an experience is putting them in the past, and the past they want is the imperfetto:

\begin{quote}
\emph{Prima dell'esame \textbf{ero nervoso}. Poi \textbf{ero felice}.} — Before the exam I was nervous. Then I was happy.
\end{quote}

A feeling has no edge. It was going on, like the weather, so in a story it is usually scene: \textbf{ero}. \emph{Sono stato felice} exists too, for a feeling summed up as one finished stretch: \emph{sono stato felice a Roma}, the whole stay.

And the adjective agrees with whoever felt it. \emph{Contento} is a regular four-ending adjective: \textbf{contento}, \textbf{contenta}, \textbf{contenti}, \textbf{contente}. \emph{Felice} and \emph{triste} end in \emph{-e} and have only two: \textbf{felice} for one, \textbf{felici} for more, man or woman.

\begin{quote}
\emph{Eravamo contenti.} — We were happy. \emph{Mia madre era felice.} — My mother was happy.
\end{quote}
\end{grammarlens}

\begin{grammarlens}[title={Grammar Lens: contento di}]
To say \textbf{what} you were happy about, add \textbf{di}:

\begin{quote}
\emph{Ero contento \textbf{del} viaggio.} — I was pleased with the journey. \emph{Sono felice \textbf{di} essere qui.} — I'm happy to be here.
\end{quote}

\emph{di} + \emph{il} is \textbf{del}, the same welding as \emph{a} + \emph{il} in \emph{al}. Before an infinitive, \emph{di} stays as it is.
\end{grammarlens}

\begin{cousinweb}
\textbf{contento} comes from Latin \emph{contentus}, \textquotedblleft{}satisfied\textquotedblright{}, the participle of \emph{continēre}, \textquotedblleft{}to hold together, to contain\textquotedblright{}, itself \emph{con-} + \emph{tenēre}, \textquotedblleft{}to hold\textquotedblright{}. Lewis \& Short, under \emph{contineo}, tie \emph{contentus} to the verb's sense \textquotedblleft{}to hold back, to keep in check\textquotedblright{}: to be content is to hold yourself in, wanting nothing further. English kept both senses apart: the \textbf{content} of a box, and a \textbf{content} person. EtymDB 2.1 links the Italian word to \emph{contentus}, though it also mixes in a different Latin word, \emph{contemptus}, \textquotedblleft{}scorn\textquotedblright{}, which is not its source.
\end{cousinweb}

\subsection*{Your turn}
Say these aloud:

\begin{itemize}
\raggedright
  \item \emph{Prima dell'esame ero nervoso. Poi ero felice.}
  \item the four endings — \emph{contento, contenta, contenti, contente}
  \item \emph{Ero contento del viaggio.}
  \item a feeling of yours, last week — \emph{La settimana scorsa ero …}
\end{itemize}

\begin{tcolorbox}[breakable,colback=teal!4,colframe=teal!35!black,title={Before you move on}]
Why is \emph{ero felice} the usual choice in a story? (A feeling has \textbf{no edge}: it is scene.) How many forms has \emph{felice}? (\textbf{Two} — \emph{felice}, \emph{felici}.) Say \textquotedblleft{}I was pleased with the journey\textquotedblright{}. (\emph{\textbf{Ero contento del viaggio.}}) What did Latin \emph{contentus} mean? (\textbf{Satisfied} — held in, wanting nothing more.)
\end{tcolorbox}
\section[spero di]{spero di — saying what you are reaching for}

\label{lesson:IT-C269-spero-di}



\begin{quote}
Say \textquotedblleft{}I'm happy to be here\textquotedblright{}. (\emph{Sono felice di essere qui.}) Then \textquotedblleft{}I would like to travel\textquotedblright{}. (\emph{Vorrei viaggiare.})

An experience has a past and a future. Last lesson was how it felt; this one is what you want next.
\end{quote}

\subsection*{You'll want to know: spero di}
\textbf{spero di} + infinitive — \textquotedblleft{}I hope to\textquotedblright{}.

\begin{quote}
\textbf{Spero di} viaggiare. — I hope to travel. \textbf{Spero di} imparare bene l'italiano. — I hope to learn Italian well.
\end{quote}

You have \textbf{sperare}, \textquotedblleft{}to hope\textquotedblright{}; \textbf{spero} is its \textquotedblleft{}I\textquotedblright{} form. The \textbf{di} is the same small word as in \emph{contento di} and \emph{credo di sì}: it hooks the infinitive on, and Italian will not leave it out. English says \emph{I hope to}; Italian says \emph{spero di}.

The \emph{di} works when the hoper and the doer are the same person. Hoping something for \textbf{someone else} takes \emph{spero che} and the \emph{congiuntivo}, which is a later chapter's work.

\begin{grammarlens}[title={Grammar Lens: il mio sogno è}]
An ambition is bigger than a hope, and it has its own frame:

\begin{quote}
\textbf{Il mio sogno è} vivere al mare. — My dream is to live by the sea. \textbf{Il mio sogno è} studiare all'università. — My dream is to study at university.
\end{quote}

\textbf{il sogno}, \textquotedblleft{}the dream\textquotedblright{}, you have; the article stays in front of \emph{mio}, as it has throughout this book. After \emph{è} the infinitive follows bare, with no \emph{di}: the dream \textbf{is} the doing.

Three strengths, then, for the future: \emph{vorrei viaggiare}, I'd like to; \emph{spero di viaggiare}, I hope to; \emph{il mio sogno è viaggiare}, it is my dream.
\end{grammarlens}

\begin{cousinweb}
\textbf{sperare} is inherited from Latin \emph{spērāre}, \textquotedblleft{}to hope, to expect\textquotedblright{} (EtymDB 2.1; Lewis \& Short, under \emph{spero}). English \textbf{desperate} is the same verb with \emph{dē-} in front, \textquotedblleft{}hope taken away\textquotedblright{}, and \textbf{la speranza}, \textquotedblleft{}hope\textquotedblright{}, which you met in chapter one hundred and eighty-six, is its noun.
\end{cousinweb}

\subsection*{Your turn}
Say these aloud:

\begin{itemize}
\raggedright
  \item \emph{Spero di viaggiare.}
  \item \emph{Spero di imparare bene l'italiano.}
  \item \emph{Il mio sogno è vivere al mare.}
  \item the three strengths — \emph{vorrei}, \emph{spero di}, \emph{il mio sogno è} — with one verb of your own
\end{itemize}

\begin{tcolorbox}[breakable,colback=teal!4,colframe=teal!35!black,title={Before you move on}]
What goes between \emph{spero} and the infinitive? (\emph{\textbf{Di.}}) And after \emph{il mio sogno è}? (\textbf{Nothing} — the infinitive follows bare.) Which English word is \emph{spērāre} with \emph{dē-} in front? (\textbf{Desperate.})
\end{tcolorbox}
\section[(review)]{Review — how it felt, and what comes next}

\label{lesson:IT-C269-ripasso-sentimenti}



\begin{quote}
Say \textquotedblleft{}I hope to travel\textquotedblright{}. (\emph{Spero di viaggiare.})
\end{quote}

\begin{grammarlens}[title={What you've built}]
An experience told well has two ends. The feeling looks back, in the imperfetto; the hope looks forward, with \emph{spero di} or \emph{il mio sogno è}.

\begin{quote}
\emph{Il viaggio era lungo. Ho perso la coincidenza, quindi ero arrabbiato. Ma alla fine ero contento del viaggio. Spero di tornare.}
\end{quote}

Follow the feelings. \emph{Ero arrabbiato} came from an event, \emph{ho perso}, and \emph{quindi} joined them; \emph{ero contento} closed the story; \emph{spero di tornare} opened the next one.
\end{grammarlens}

\subsection*{Your turn}
Say these aloud:

\begin{itemize}
\raggedright
  \item \emph{Ho perso la coincidenza, quindi ero arrabbiato.}
  \item \emph{Ma alla fine ero contento del viaggio.}
  \item \emph{Spero di tornare.}
  \item a journey of your own: one feeling, then a \emph{poi}, then a second feeling
  \item your dream — \emph{Il mio sogno è …}
\end{itemize}

\begin{tcolorbox}[breakable,colback=teal!4,colframe=teal!35!black,title={Before you move on}]
Which past do the feelings take? (The \textbf{imperfetto}: \emph{ero arrabbiato}, \emph{ero contento}.) Which one is an event? (\emph{\textbf{Ho perso.}}) What hooks an infinitive onto \emph{spero}? (\emph{\textbf{Di.}})
\end{tcolorbox}
